\documentclass[12pt]{article}
\usepackage[T1]{fontenc}
\usepackage[utf8]{inputenc}
\usepackage{lmodern}
\usepackage{textcomp}
\usepackage{enumitem}
\usepackage{geometry}
\geometry{margin=1in}
\begin{document}
\begin{center}
Answer Key - Health Sciences - K-12 Level Assessment 24 \
\small (Answers are ordered exactly as in the question paper.)
\end{center}

\section*{Multiple Choice}
\begin{enumerate}[label=\arabic*.]
\item \textbf{Answer:} D
\\ \textit{Options:} A) Decreased calcium excretion in the urine; B) Increased risk of fracture; C) Decreased dietary calcium absorption; D) Negative calcium balance and bone mineral loss
\\ \textit{Note:} Higher dietary sodium intake can lead to increased calcium excretion in the urine, resulting in a negative calcium balance and potential bone mineral loss over time.
\item \textbf{Answer:} D
\\ \textit{Options:} A) O$_{2}$ consumption alone; B) CO 2 production alone; C) Breath collection; D) Total urinary nitrogen excretion alone
\\ \textit{Note:} Total urinary nitrogen excretion alone is a reliable measure of protein oxidation because the nitrogen from oxidized protein is primarily excreted through urine, allowing for the estimation of protein catabolism in the body.
\item \textbf{Answer:} A
\\ \textit{Options:} A) \textless{}10\% of total daily energy; B) \textless{}12\%; C) \textless{}15\%; D) \textless{}16\%
\\ \textit{Note:} The recommendation for patients with diabetes to limit saturated fat to less than 10\% of total daily energy is based on evidence linking high saturated fat intake to increased risk of cardiovascular disease and poor glycemic control, which are particularly concerning for individuals managing diabetes.
\item \textbf{Answer:} C
\\ \textit{Options:} A) Normal energy intake, normal body weight and hyperthyroidism; B) Obesity, excess energy intake, normal growth and hypoinsulinaemia; C) Obesity, abnormal growth, hypothyroidism, hyperinsulinaemia; D) Underweight, abnormal growth, hypothyroidism, hyperinsulinaemia
\\ \textit{Note:} The correct answer is based on the fact that impaired leptin secretion disrupts appetite regulation and energy balance, leading to obesity, abnormal growth due to altered metabolic processes, hypothyroidism from reduced energy expenditure, and hyperinsulinaemia due to increased insulin resistance.
\item \textbf{Answer:} A
\\ \textit{Options:} A) DNMT 1; B) DNMT 3 a; C) DNMT 3 b; D) DNMT 3 L
\\ \textit{Note:} DNMT 1 is the enzyme specifically responsible for maintaining DNA methylation patterns by transferring methyl groups from the parental DNA strand to the newly synthesized daughter strand during DNA replication.
\end{enumerate}

\section*{True / False}
\begin{enumerate}[label=\arabic*.]
\setcounter{enumi}{5}
\item \textbf{Answer:} True
\\ \textit{Options:} A) True; B) False
\\ \textit{Note:} Starch is primarily composed of amylose, which is a polysaccharide made up of glucose monomers linked by alpha 1,4-glycosidic bonds, confirming the statement as true.
\item \textbf{Answer:} False
\\ \textit{Options:} A) True; B) False
\\ \textit{Note:} C$_{18}$ and C$_{22}$ polyunsaturated fatty acids are not the primary precursors for eicosanoid synthesis; instead, C$_{20}$ polyunsaturated fatty acids, particularly arachidonic acid (AA), play a crucial role in the production of eicosanoids.
\item \textbf{Answer:} False
\\ \textit{Options:} A) True; B) False
\\ \textit{Note:} Supplementation with EPA and DHA during pregnancy is primarily beneficial for overall brain development and visual acuity in infants, rather than specifically improving memory.
\item \textbf{Answer:} False
\\ \textit{Options:} A) True; B) False
\\ \textit{Note:} Saturated and trans unsaturated fatty acids are known to raise plasma LDL cholesterol levels, which can increase the risk of heart disease, rather than improve them.
\item \textbf{Answer:} False
\\ \textit{Options:} A) True; B) False
\\ \textit{Note:} The statement is false because the basic functional unit of skeletal muscle tissue is the sarcomere, which is composed of muscle fibers and is responsible for muscle contraction.
\end{enumerate}

\section*{Long Form Response}
\begin{enumerate}[label=\arabic*.]
\setcounter{enumi}{10}
\item \textbf{Answer:} In the fed state, when we consume food, our body experiences increased insulin secretion, which plays a crucial role in promoting tissue protein synthesis. Insulin helps transport amino acids into cells, enabling them to build proteins that are essential for growth, repair, and overall maintenance of body tissues. This process is significant for overall health because adequate protein synthesis supports muscle development, immune function, and recovery from injuries. Additionally, proper protein synthesis helps maintain a balanced metabolism and can prevent conditions like malnutrition. Therefore, increased insulin secretion after eating is vital for ensuring that our bodies have the resources they need to function optimally.
\\ \textit{Note:} correct because it emphasizes that increased insulin secretion in the fed state enhances the transport of amino acids into cells, thereby promoting protein synthesis, which is vital for growth, repair, and overall health.
\item \textbf{Answer:} High doses of ß-carotene supplements may increase the risk of developing lung cancer, especially among smokers. Research has shown that while ß-carotene is a type of antioxidant found in fruits and vegetables, taking it in high doses through supplements can have harmful effects. One possible mechanism is that excessive amounts of ß-carotene can lead to oxidative stress, which may promote cancer cell growth. Studies, including a large clinical trial, indicated that smokers who took ß-carotene supplements had a higher incidence of lung cancer compared to those who did not. Therefore, it is important to be cautious with high-dose supplements and to focus on getting nutrients from a balanced diet instead.
\\ \textit{Note:} correct because it identifies that high doses of ß-carotene supplements can lead to increased oxidative stress, particularly in smokers, which has been supported by research indicating a higher risk of lung cancer in this population.
\end{enumerate}

\vfill
\noindent\textit{End of Answer Key}
\end{document}